\section{A joint chart for the shared point and its translate}

This section starts from the 423-coordinate construction.  Work over
$K=\mathbb Q(\rho)$, $\rho^2=2$, and put $D=(3-\rho)^2$.
Write the fixed point as $P=(x,hj)$, where $x=h^2$, $j^2=x^2-D$.
Here $h,x,j$ are the fixed nonzero constants already used in the source.
The shared point $H$ and $V=H-P$ must both be nonidentity.

Instead of the previous coordinates $g,l,f,\tau,F,M,k_H$, use
\[
 g,\ s,\ f,\ F,\ k_H,\ k_V
\]
and the five quadratic equations over $K$
\begin{align}
 f&=g^2, & fF&=s^2,\label{r4:eq:joint-square}\\
 Ds^2+x(f+F)+2js&=1,\label{r4:eq:joint-curve}\\
 hs k_H&=1-xf-js, & hs k_V&=xF+js-1.\label{r4:eq:joint-ratios}
\end{align}
The constants $h,x,j,D$ do not add coordinates.  Each $K$-coordinate
consists of two rational coordinates, so replacing seven by six saves two.

\subsection{All denominators are forced nonzero}
If $f=0$, then $s^2=fF=0$, and the first ratio equation says $0=1$.
If $F=0$, then $s=0$, and the second ratio equation says $0=-1$.
Thus $f,F,s,g$ are nonzero, without adding guards.  The implemented
quotients $F=s^2/f$ and the two divisions by $hs$ are consequently
well-defined on every solution of the full joint system.

\subsection{Exact recovery of the old complete chart}
Set
\[
 \tau=s/g,\qquad l=gk_H,\qquad M=j+hk_V.
\]
Then $\tau^2=F$.  With $B=x+Df$, equations
\eqref{r4:eq:joint-square}--\eqref{r4:eq:joint-ratios} give
\[
 B\tau=hl-gj,\qquad
 \tau M=-g(x+DF).
\]
For the first identity, multiply by $gs$: the difference reduces to
$-x(fF-s^2)-f(hsk_H-1+xf+js)+f(Ds^2+x(f+F)+2js-1)$.
For the second, multiply by $g$ and use
\[
 s(j+hk_V)+f(x+DF)
 =(hsk_V-xF-js+1)+(Ds^2+x(f+F)+2js-1)+D(fF-s^2).
\]

The point-chart equation follows as well.  Put
$A=hl-gj=B\tau$.  Then
\[
 A(hl+gj)=B\bigl(BF+2js\bigr)=B(1-xf),
\]
where the middle equality uses $g\tau=s$ and the last uses
\eqref{r4:eq:joint-curve}.  On the fixed root chart,
\[
 (hl-gj)(hl+gj)-(1-xg^2)(x+Dg^2)
 =x(l^2-1+Dg^4)-g^2(j^2-x^2+D).
\]
Since $x\ne0$, this gives $l^2=1-Df^2$.
Consequently $H=(1/f,k_H/f)$ is exactly the old reciprocal-chart point,
and the recovered equations are exactly the preceding complete
subtraction equations for $H-P$.  They include $H=-P$ and exclude
$H=P$, as required.  No ordinate sign or extra group component is introduced.

Conversely, from the preceding valid shared-point chart set $s=g\tau$.
The same identities give \eqref{r4:eq:joint-curve} and both ratio equations.
The two maps are inverse on their forced nonzero domains.  This change is
therefore a genuine rational coordinate equivalence, preserving the
actual $H,V,f_H,f_V,k_H,k_V$ functions.

For completeness, $k_H$ and $k_V$ cannot vanish.  A reciprocal square
chart with $k=0$ would give $Dg^4=1$, hence
$7N_{K/\mathbb Q}(g)^2=1$, impossible over $\mathbb Q$.
The old nonzero logarithm denominators are thus preserved.  There is no
new local or global arithmetic input.

\section{Fuse the scalar reciprocal difference}

Let $S=(x_S,h_Sj_S)$ retain its root chart, so $x_S=h_S^2\ne0$.
The previous scalar difference $R=S-P$ used
\[
 B\tau=h j_S-h_Sj,\qquad F=\tau^2,\qquad
 w=h_S\tau,\qquad wM=-(x+DF),\qquad B=xx_S+D.
\]
Replace its first three equations by
\[
 \boxed{Bw=h h_Sj_S-x_Sj,\qquad x_SF=w^2,}
\]
and retain the equation for $M$ and $k_R=(M-j)/h$.
The new numerator is quadratic; $B$ is affine and totally positive.
This removes the two rational coordinates of $\tau$.

The forward map is $w=h_S\tau$.  Conversely define $\tau=w/h_S$.
Then $x_S=h_S^2$ recovers the old first equation and $F=\tau^2$.
If $w=0$, the new reciprocal equation gives $F=0$, and the retained
$M$ equation would force $x=0$.  Thus $w\ne0$ is forced before dividing.
All previous scalar point functions are unchanged.  Tangency $S=-P$
remains allowed, and the identity output $S=P$ remains excluded.

The two coordinate changes above give 419 variables and 387 residuals
in isolation, compared with 423 and 391.  These are compiled counts.
They retain quadratic residuals when the fixed coefficient $S_0$ is
subsequently replaced by a constant-chain coordinate.

\section{Inverse-root tests and the shared primary depth}

At a consumed unramified place, normalize the valuation by $v(p)=1$.
For $f=g^2$ one has
\[
 v(f)\ge1\quad\Longleftrightarrow\quad v(g)\ge1,
\]
including $g=0$ with valuation $+\infty$.  Therefore the order-one
primary tests on $f(U_i)$ can be replaced by tests on $g_i$ at the same
order.  Their squares are the already materialized coordinates $f(U_i)$.
The alternate tests may be left as $f(V_i)$; this is compatible with the
joint shared chart, which does not materialize a separate inverse root
for $V_0$.

Replacing the shared order-40 test on $f(H)$ by an order-40 test on
$g_H$ is a stronger condition: it gives primary formal depth at least
40, rather than 20.  Soundness uses only the old lower bound, so all its
implications continue to hold.

For completeness, use the actual CRT passage in
\texttt{docs/quartic-523-540/addition.tex}, in the subsections
``A stronger choice within the existing prime-pattern argument'' and
``Completeness with deeper primary points.''  The roots, scales and
continuity modulus $L$ are fixed before the primes.  Let
\[
 M_\pm=\operatorname{lcm}_p(d_\pm(p)/4),\qquad \Pi=\prod_p p.
\]
The source supplies
\[
 \gcd(M_-,M_+)=1,\quad \gcd(M_-,L)=16,\quad
 \gcd(M_+,L)=1,\quad \gcd(\Pi,M_+L)=1,
\]
with $64\mid h_0$.  Hence the stronger congruences
\[
 h\equiv0\pmod{\Pi^{39}M_-},\qquad
 h\equiv1\pmod{M_+},\qquad h\equiv h_0\pmod L
\]
are compatible by the same CRT argument.  The inherited real progression
argument applies unchanged.  The formal-logarithm isometry raises the
primary depth from at least one to at least forty.  Thus $H=[h]P$ and
$U_i=[hn_i]P$ meet all new primary tests.  The alternate points retain
their old depth-one requirements; no contradictory divisibility is
imposed on $h-1$.

In the detector products the new factors use $g_i^2=f(U_i)$.
Thus their dependence on the actual point remains rational and
independent of the sign of the square-root chart.  The source's finite-prime
and real continuity argument can retain the same prepared point after
enlarging its fixed modulus or bounds as needed.  The separate detector
review checks the unchanged general pole estimates for the chosen test list.

This stronger shared test is a witness-space restriction with a checked
completeness construction; it is not asserted to be a bijection on all
old witnesses.

\section{Only the input comparison needs a degree-two alternate logarithm}

The source's height contradiction for the Lucas residual uses the three
root indices $a,t,b$ and modulus $p^{40}$.  Its contradiction for
$n-\beta$ uses only modulus $p^{12}$; these are the explicit inequalities
in \texttt{addition.tex}, the subsection ``Soundness'' of
``Odd labels, scalar bounds, and completeness.''

Keep $G_9$ for $V_0,V_2,V_3,V_4$, and use
\[
 G_2(w)=1-\frac25w-\frac{2}{15}w^2
\]
for $V_1$, the translate corresponding to the input root $A=[n]P$.
Change only the cross test $L_1-L_1'$ to order 12; the beta test may
retain order 40.  The three other cross tests and the Lucas test retain
order 40.

If $Q$ is a formal point of depth $r\ge1$, then
$v(Df(Q)^2)=4r$.  The source coefficient bound
$v(a_j)\ge-(4j+1)/(p-1)$, together with integrality when $4j+1<p$,
shows that the tail after $G_2$ has valuation at least $12r$ at the
consumed primes $p>76176$.  Thus
\[
 \frac{\log_E(Q)}{z(Q)G_2(Df(Q)^2)}\in1+p^{12}\mathcal O.
\]
All denominators are nonzero: on the finite real chart $0<Df(Q)^2<1$,
and $G_2(w)\ge G_2(1)=7/15>0$.

The mixed ratio
\[
 L_1'=\overline{\frac{\ell_2(V_1)}{\ell_9(V_0)}}
\]
therefore has relative precision 12 against
$r_1'=(u_1-n)/(h-1)$ and has the same valuation.  The primary ratio has
the same valuation as $r_1=u_1/h$ on the old formal domain.  One of
$h,h-1$ is a $p$-unit, so one ratio is integral.  The order-12 cross
test transfers integrality to both.  Integrality of $L_1$ raises the
depth of $U_1$ to that of $H$, so the primary relative error is more than
sufficient.  Consequently
\[
 L_1\equiv r_1\equiv r_1'\pmod{p^{12}},
 \qquad n\equiv L_1\equiv\beta\pmod{p^{12}}.
\]
If $U_1=O$, the same conclusion follows from $L_1=r_1=0$ and
$n=-(h-1)r_1'$.  At inert places conjugation preserves the valuation.
All other comparison conclusions retain precision 40.  The source's
two height contradictions therefore retain precisely their used hypotheses.

For completeness, the old choices $u_i=hn_i$ make both full logarithm
ratios equal $n_i$.  The shorter alternate logarithm meets the order-12
cross test, while the primary ratio still meets the order-40 beta test.
No extra alternate depth or new prime selection is required.

Implementation stores the same scaled coordinate $\eta_1=\ell_2(V_1)/h_A$,
with
\[
 (M_1-j_A)\eta_1=-G_2(Df(V_1)^2).
\]
Its final quadratic stage inlines exactly as before.  The isolated
logarithm change saves six rational coordinates; the full saving must
also account for the detector's regrouping of orders 1,12,40.

\section{Search results that were not counted}

Raising every alternate reciprocal-abscissa order from 1 to 3 would
permit $G_4$, but the canonical completeness argument generally requires
$p\mid h-1$ to obtain that extra alternate depth, whereas the primary
depth already imposes $p\mid h$.  The fixed point's exceptional higher
formal depth at a selected prime is not supplied by the prime-pattern
theorem.  This route is therefore not accepted.

A generic Pad\'e numerator/denominator of degrees 4 and 5 also does not
give a coordinate reduction by itself.  Its denominator requires an
additional product with the logarithm denominator.  The implemented
$G_9$ block already reaches degree nine with three materialized intermediate
$K$-coordinates followed by a quadratic quotient row.  No Pad\'e saving
is asserted without a separate circuit and domain proof.

All global arithmetic inputs remain inherited.  The exact checks supplied
with this note establish their specified polynomial/source identities,
not the truth of those inherited results.
